% !TeX root = ../../Book2.tex
\OptionalSection{一阶非线性方程}{b2:sec:0704}

在线性输运中，先由已知速度场确定轨道，再沿轨道计算数值。若速度依赖未知函数，这两步便相互耦合。沿特征线同时求解位置与函数值，可以构造局部光滑解；显式例子则会显示梯度怎样在有限时间内失去控制。

\subsection{拟线性方程}
\label{b2:subsec:070401}

考虑已在\secref{b2:sec:0104}定义的拟线性类型
\[
 u_t+b(x,t,u)\cdot\nabla u=f(x,t,u),\quad u(x,0)=g(x).
\]
这里 \(d\geq1\) 为整数，未知函数及数据 \(g\) 取实值，速度 \(b\) 取值于 \(\mathbb R^d\)，源项 \(f\) 取实值。设 \(b,f\) 在相关变量的开集上为 \(C^1\)，\(g\in C^1\)。把未知的函数值也看成一个坐标，解联立系统
\begin{equation}\label{b2:eq:07-quasilinear-system}
 \begin{aligned}
 X'&=b(X,t,Z),&X(0,a)&=a,\\
 Z'&=f(X,t,Z),&Z(0,a)&=g(a).
 \end{aligned}
\end{equation}
将时间也作为坐标，补上 \(t'=1\)，便可对向量 \((X,t,Z)\) 应用\cref{b2:thm:07-flow-derivative}，得到其对 \((t,a)\) 的 \(C^1\) 依赖。时刻零有 \(D_aX=I_d\)，故映射 \((a,t)\mapsto(X(t,a),t)\) 在每个 \((a_0,0)\) 附近的导数可逆。

\begin{proposition}\label{b2:prop:07-quasilinear-local}
设 \(d\geq1\) 为整数，\(V\subset\mathbb R^d\) 为开集，\(g\in C^1(V;\mathbb R)\)，\(a_0\in V\)，\(b,f\) 在包含 \((a_0,0,g(a_0))\) 的 \(\mathbb R^d\times\mathbb R\times\mathbb R\) 开集上属于 \(C^1\)，值域分别为 \(\mathbb R^d\)、\(\mathbb R\)。则在 \((a_0,0)\) 的某个时空邻域中，初值问题 \(u_t+b(x,t,u)\cdot\nabla u=f(x,t,u)\)、\(u(x,0)=g(x)\) 有实值 \(C^1\) 解，且该解局部唯一；任何满足相同初值的另一实值 \(C^1\) 解，在适当缩小的共同邻域内与它相等。
\end{proposition}
\begin{proof}
由\cref{b2:thm:01-inverse}，缩小邻域后可反解 \(a=A(x,t)\)。令 \(u(x,t)=Z(t,A(x,t))\)，则
\(u(X(t,a),t)=Z(t,a)\)。对此式求时间导数并代入\eqref{b2:eq:07-quasilinear-system}，得到偏微分方程。初始时 \(A(x,0)=x\)，所以 \(u(x,0)=g(x)\)。若另有经典解 \(v\)，沿满足 \(X'=b(X,t,v(X,t))\) 的曲线令 \(Z=v(X,t)\)，链式法则表明 \((X,Z)\) 满足同一联立系统和同一初值。常微分方程唯一性使二者与构造的 \((X,Z)\) 相同，再由局部反演得到 \(v=u\)。
\end{proof}

这一结论的邻域可以依赖 \(a_0\)。即使联立系统一直有解，其空间投影 \(a\mapsto X(t,a)\) 也可能先失去可逆性。下面的例子把这种障碍显示为一个分母趋零的计算。

\subsection{特征线交叉}
\label{b2:subsec:070402}

\term[burgers-fangcheng]{无黏 Burgers 方程}[inviscid Burgers equation]为
\[
 u_t+u u_x=0.
\]
取初值 \(u(x,0)=g(x)\)，其中 \(g\in C^1(\mathbb R;\mathbb R)\) 有界，考察 \(t\geq0\)。沿特征线 \(Z'=0\)，所以 \(Z(t,a)=g(a)\)、\(X(t,a)=a+tg(a)\)。在满足 \(1+tg'(a)>0\) 对所有 \(a\in\mathbb R\) 成立的时间内，\(X(t,\cdot)\) 严格递增，又因 \(|X(t,a)-a|\leq t\|g\|_\infty\) 而满射。它是全实轴的 \(C^1\) 微分同胚，局部逆随时空变量 \(C^1\) 变化，解由 \(u(X(t,a),t)=g(a)\) 给出，并满足
\begin{equation}\label{b2:eq:07-burgers-gradient}
 u_x(X(t,a),t)=\frac{g'(a)}{1+tg'(a)}.
\end{equation}
若 \(g'\geq0\)，上述构造对所有正时间有效；若 \(g'\) 的负最小值 \(m\) 在某点取到，则在 \(T_*=-1/m\) 之前有效，而该点轨道上的梯度在 \(t\uparrow T_*\) 时趋于负无穷。

\ProseParagraphBreak

具体取 \(g(a)=-\tanh a\)，有 \(m=-1\)、\(T_*=1\)。在 \(0\leq t<1\) 时，函数值始终处于 \((-1,1)\)，但原点处梯度为 \(-1/(1-t)\)。失效来自投影的导数 \(\partial_aX(t,0)=1-t\) 趋于零，即相邻初始标签之间的无穷小空间间隔被压缩。在 \(t=1\) 时，\(a\mapsto a-\tanh a\) 仍严格递增，但在原点的导数为零，逆映射已不再为 \(C^1\)。因此首次梯度爆破并不要求两条不同标签的投影恰在这一时刻相交。当 \(t>1\) 时，映射 \(a-t\tanh a\) 在原点附近下降，在两端上升，才失去单射性。两条不同初始标签的轨道可以投影到同一个时空点，而所携带的值不同，不可能成为单值经典解。\Cref{b2:fig:07-characteristic-compression}区分了光滑阶段与失效后的轨道投影。

\ProseParagraphBreak

这里交叉的是 \((x,t)\) 平面中的投影。在扩充状态 \((X,Z)\) 中，特征系统为 \(X'=Z\)、\(Z'=0\)，其流是 \((a,z_0)\mapsto(a+tz_0,z_0)\)，逆为 \((x,z)\mapsto(x-tz,z)\)。所以完整状态的流始终单射，常微分方程的唯一性并未失效；失效的是丢去 \(Z\) 坐标后，用位置反解初始标签的步骤。
% 自绘：g(a)=-tanh(a) 的特征投影；虚线仅表示光滑失效后的轨道投影。
\begin{figure}[htbp]
\centering
\begin{tikzpicture}[x=2.3cm,y=2.2cm,>=stealth]
 \draw[->] (-1.35,0)--(1.35,0) node[right]{\(x\)};
 \draw[->] (-1.3,-0.05)--(-1.3,1.6) node[above]{\(t\)};
 \draw[gray,densely dotted] (-1.25,1)--(1.25,1);
 \node[left] at (-1.3,1){\(1\)};
 \foreach \a/\v in {-1/-0.761594,-0.75/-0.635149,-0.5/-0.462117,-0.25/-0.244919,0/0,0.25/0.244919,0.5/0.462117,0.75/0.635149,1/0.761594}{
   \draw[AnalysisBlue,thick] (\a,0)--({\a-\v},1);
   \draw[AnalysisBlue,dashed] ({\a-\v},1)--({\a-1.45*\v},1.45);
 }
 \foreach \x in {-1,0,1}{\draw (\x,0.025)--(\x,-0.025) node[below]{\(\x\)};}
 \node[anchor=west,align=left] at (0.68,1.28){\small 轨道投影};
 \node[anchor=west] at (0.68,0.72){\small 光滑阶段};
\end{tikzpicture}
\caption[特征线的压缩与投影交叉]{特征线的压缩与投影交叉。初值为 \(g(a)=-\tanh a\)，轨道投影为 \(x=a-t\tanh a\)。在 \(t=1\) 首次发生梯度爆破；虚线延长只展示投影的交叉，不表示单值经典解的延续。}
\label{b2:fig:07-characteristic-compression}
\end{figure}


\subsection{激波形成的预告}
\label{b2:subsec:070403}

在经典解存在时，Burgers 方程等价于
\(u_t+\partial_x(u^2/2)=0\)。对于有跳跃的函数，后一种写法仍可通过分布导数解释。设一般通量 \(F\in C^1(\mathbb R;\mathbb R)\)，\(I\subset\mathbb R\) 为开时间区间，\(s\in C^1(I;\mathbb R)\)。令实值函数在 \(x<s(t)\) 一侧恒为常数 \(u_-\)，在 \(x>s(t)\) 一侧恒为常数 \(u_+\)；曲线上的取值不影响分布。对 \(\varphi\in C_c^\infty(\mathbb R\times I)\)，把分布 \(u_t+\partial_xF(u)\) 的测试积分在曲线两侧分别分部积分，得到界面项
\[
 \int \bigl(F(u_+)-F(u_-)-s'(t)(u_+-u_-)\bigr)
 \varphi(s(t),t)\,\mathrm dt.
\]
它对所有紧支集光滑测试函数为零，当且仅当
\[
 s'(t)(u_+-u_-)=F(u_+)-F(u_-).
\]
这称为\term[rankine-hugoniot]{Rankine--Hugoniot 条件}[Rankine--Hugoniot condition]。对不同的 Burgers 左右状态，速度为 \((u_-+u_+)/2\)；这样的移动跳跃称为\term[jibo]{激波}[shock wave]。

\ProseParagraphBreak

界面恒等式只给出弱方程的相容条件，没有证明某个跳跃是光滑解失效后的唯一延续。物理上允许的压缩跳跃与不允许的展开跳跃，都可能满足同一恒等式。熵条件用来区分这些跳跃；\secref{b2:sec:0705}将从两常数初值出发，计算激波与稀疏波，并解释相应的选择。

\subsection{Hamilton--Jacobi 方程}
\label{b2:subsec:070404}

另一类重要的一阶非线性方程是\term[hamilton-jacobi-fangcheng]{Hamilton--Jacobi 方程}[Hamilton--Jacobi equation]
\[
 u_t+H(x,\nabla u)=0.
\]
这里 \(d\geq1\) 为整数，\(H\) 与未知函数 \(u\) 都取实值。设 \(V\subset\mathbb R^d\) 开、\(a_0\in V\)，\(g\in C^3(V;\mathbb R)\)，\(H\) 是包含 \((a_0,\nabla g(a_0))\) 的 \(\mathbb R^d\times\mathbb R^d\) 开邻域上的实值 \(C^3\) 函数，只在初始点附近构造光滑解。取
\begin{equation}\label{b2:eq:07-hamilton-system}
 X'=H_p(X,P),\quad P'=-H_x(X,P),\quad
 Z'=P\cdot H_p(X,P)-H(X,P),
\end{equation}
初始值为 \(X(0,a)=a\)、\(P(0,a)=\nabla g(a)\)、\(Z(0,a)=g(a)\)。前两式描述位置和候选梯度的演化，第三式计算函数值。

\begin{proposition}\label{b2:prop:07-hamilton-local}
设 \(d\geq1\) 为整数，\(V\subset\mathbb R^d\) 为开集，\(g\in C^3(V;\mathbb R)\)，\(a_0\in V\)。设 \(\Omega\subset\mathbb R^d\times\mathbb R^d\) 为包含 \((a_0,\nabla g(a_0))\) 的开集，\(H\in C^3(\Omega;\mathbb R)\)。取 \(X(0,a)=a\)、\(P(0,a)=\nabla g(a)\)、\(Z(0,a)=g(a)\)，系统\eqref{b2:eq:07-hamilton-system}在 \((a_0,0)\) 的某个时空邻域中确定一个实值 \(C^2\) 解 \(u\)，满足 \(u_t+H(x,\nabla u)=0\)、\(u(x,0)=g(x)\) 及 \((x,\nabla u(x,t))\in\Omega\)。该解在实值 \(C^2\) 类中局部唯一：若另一实值 \(C^2\) 解 \(v\) 定义在 \((a_0,0)\) 附近，满足同一方程和初值，且 \((x,\nabla v(x,t))\in\Omega\)，则 \(v\) 在适当缩小的共同邻域内与 \(u\) 相等。
\end{proposition}
\begin{proof}
先核对特征系统的二阶参数依赖。令 \(N=2d+1\)，把独立初始状态写成 \(y_0=(x_0,p_0,z_0)\in\Omega\times\mathbb R\)，并设
\[
 B(x,p,z)=\bigl(H_p(x,p),-H_x(x,p),p\cdot H_p(x,p)-H(x,p)\bigr).
\]
由 \(H\in C^3\)，向量场 \(B\) 属于 \(C^2\)。记 \(\Phi_t(y_0)\) 为 \(Y'=B(Y)\)、\(Y(0)=y_0\) 的局部流。由\cref{b2:thm:07-flow-derivative}，流为 \(C^1\)，且 \(L(t,y_0)=D_{y_0}\Phi_t(y_0)\) 满足
\[
 L'=DB(\Phi_t(y_0))L,\quad L(0,y_0)=I_N.
\]
对独立的方向 \(r,s\in\mathbb R^N\)，二阶导数的候选值 \(V_{r,s}\) 由下列连续系数线性积分方程的等价微分形式确定；其存在唯一性用一阶变分方程相同的积分迭代得到：
\[
 V_{r,s}'=DB(\Phi_t(y_0))V_{r,s}
       +D^2B(\Phi_t(y_0))[Lr,Ls],
 \quad V_{r,s}(0)=0.
\]
零初值的原因是 \(\Phi_0(y_0)=y_0\) 对独立变量 \(y_0\) 为恒等映射。

为验证候选值确实是二阶导数，缩小初始邻域与时间区间，使所有流及其连接线段位于同一个紧球 \(K\subset\Omega\times\mathbb R\) 中。记 \(Y_\eta(t)=\Phi_t(y_0+\eta s)\)、\(Y(t)=\Phi_t(y_0)\)，其中 \(\eta\) 为趋于零的非零实数。一阶 \(C^1\) 依赖给出 \((Y_\eta-Y)/\eta\to Ls\)，在该紧时间区间上为一致收敛。对一阶变分方程作差，令
\[
 Q_\eta(t)=\frac{L(t,y_0+\eta s)r-L(t,y_0)r}{\eta},
\]
则
\[
 Q_\eta'=DB(Y_\eta)Q_\eta
       +\frac{DB(Y_\eta)-DB(Y)}{\eta}Lr,
 \quad Q_\eta(0)=0.
\]
中值积分将这里的算子差商写成
\[
 \frac{DB(Y_\eta)-DB(Y)}{\eta}
 =\int_0^1D^2B\bigl(Y+\theta(Y_\eta-Y)\bigr)
       \left[\frac{Y_\eta-Y}{\eta},\,\cdot\right]\,\mathrm{d}\theta.
\]
\(D^2B\) 在 \(K\) 上一致连续，所以该算子差商一致趋于 \(D^2B(Y)[Ls,\cdot]\)；同时 \(DB(Y_\eta)\to DB(Y)\) 一致。二阶微分的对称性使极限非齐次项 \(D^2B(Y)[Ls,Lr]\) 正是 \(V_{r,s}\) 方程中的那一项。把 \(Q_\eta\) 的积分方程与 \(V_{r,s}\) 的方程相减，所有系数及非齐次项均有共同上界，余项一致趋零，Gronwall 估计遂给出 \(Q_\eta\to V_{r,s}\) 一致。方程中系数与初始状态连续，作差后再用同一估计，得到 \(V_{r,s}\) 对 \((t,y_0,r,s)\) 连续。因此 \(D_{y_0}\Phi_t\) 的各方向偏导数连续，\(\Phi_t\) 对 \(y_0\) 为 \(C^2\)。再由 \(\partial_t\Phi_t=B(\Phi_t)\) 和一阶变分方程，混合二阶导数与二阶时间导数也连续，故在该共同邻域内，流对 \((t,y_0)\) 联合为 \(C^2\)。

现在才代入实际初始映射。先由连续性缩小 \(a_0\) 的邻域，使 \((a,\nabla g(a))\) 均落在 \(\Omega\) 内，并使下面的初始状态落在上述局部流的共同初始邻域中。令
\[
 y_0(a)=(a,\nabla g(a),g(a)),\quad
 Y(t,a)=\Phi_t(y_0(a)).
\]
由于 \(g\in C^3\)，映射 \(y_0\) 为 \(C^2\)。二阶链式法则给出
\[
 \begin{aligned}
 D_a^2Y(t,a)[v,w]
 &=D_{y_0}^2\Phi_t(y_0(a))[Dy_0(a)v,Dy_0(a)w]\\
 &\quad+D_{y_0}\Phi_t(y_0(a))D^2y_0(a)[v,w],
 \quad v,w\in\mathbb R^d.
 \end{aligned}
\]
特别地，初始时刻的结果是
\[
 D_a^2Y(0,a)[v,w]
 =\bigl(0,D^3g(a)[v,w],D^2g(a)[v,w]\bigr).
\]
这里 \(D^3g(a)[v,w]\) 表示向量 \(D^2(\nabla g)(a)[v,w]\)，其第 \(i\) 个分量为 \(D^2(\partial_i g)(a)[v,w]\)。这个初值一般不为零；为零的是独立状态流的 \(D_{y_0}^2\Phi_0\)。这一复合论证证明了 \(X,P,Z\) 对 \((t,a)\) 为 \(C^2\)。

从 \(D_aX(0,a)=I_d\) 出发，对 \((a,t)\mapsto(X(t,a),t)\) 应用\cref{b2:thm:01-inverse}，得到局部逆 \((x,t)\mapsto(A(x,t),t)\)。原映射为 \(C^2\)，逆微分公式说明逆映射的导数为 \(C^1\)，故逆也为 \(C^2\)。令 \(u(x,t)=Z(t,A(x,t))\)，便有 \(u\in C^2\) 及 \(u(x,0)=g(x)\)。关键是验证 \(P=\nabla u\)，不能把候选梯度直接当成真实梯度。对每个初始坐标 \(a_j\)，记
\(W_j=\partial_{a_j}Z-P\cdot\partial_{a_j}X\)。对系统求导，利用 \(X'=H_p\)、\(P'=-H_x\)，得到
\[
 \begin{aligned}
 \partial_{a_j}Z'
 &=P\cdot\partial_{a_j}(H_p)-H_x\cdot\partial_{a_j}X,\\
 (P\cdot\partial_{a_j}X)'
 &=-H_x\cdot\partial_{a_j}X+P\cdot\partial_{a_j}(H_p).
 \end{aligned}
\]
故 \(W_j'=0\)，初值又为 \(\partial_jg-\partial_jg=0\)。因此 \(D_aZ=P^{\mathsf T}D_aX\)，由 \(D_aX\) 可逆得到 \(\nabla_xu(X,t)=P\)。最后沿轨道有
\(u_t+\nabla u\cdot X'=Z'=P\cdot H_p-H\)，约去相同项得 \(u_t+H(x,\nabla u)=0\)。

最后证明局部唯一性。设 \(v\) 是在 \((a_0,0)\) 附近满足同一初值问题的另一实值 \(C^2\) 解，且 \((x,\nabla v(x,t))\in\Omega\)。对它的方程求空间梯度，得到
\[
 \partial_t\nabla v+D_x^2v\,H_p(x,\nabla v)+H_x(x,\nabla v)=0.
\]
速度 \((x,t)\mapsto H_p(x,\nabla v(x,t))\) 为 \(C^1\)，故从每个邻近标签 \(a\) 出发可以解
\[
 \widetilde X'=H_p(\widetilde X,\nabla v(\widetilde X,t)),
 \quad \widetilde X(0,a)=a.
\]
缩小邻域使这些短轨道始终留在 \(v\) 的定义域内，并令 \(\widetilde P(t,a)=\nabla v(\widetilde X(t,a),t)\)、\(\widetilde Z(t,a)=v(\widetilde X(t,a),t)\)。链式法则与梯度方程分别给出
\[
 \widetilde P'=-H_x(\widetilde X,\widetilde P),\quad
 \widetilde Z'=\widetilde P\cdot H_p(\widetilde X,\widetilde P)
                    -H(\widetilde X,\widetilde P).
\]
其初始状态同为 \((a,\nabla g(a),g(a))\)。向量场 \(B\) 的常微分方程唯一性使 \((\widetilde X,\widetilde P,\widetilde Z)=(X,P,Z)\)。因此 \(v(X(t,a),t)=Z(t,a)=u(X(t,a),t)\)，局部投影的逆坐标又覆盖了一个时空邻域，遂得到该邻域内 \(v=u\)。
\end{proof}

取 \(H(p)=|p|^2/2\) 时，\(P=\nabla g(a)\)、\(X=a+t\nabla g(a)\)、\(Z=g(a)+t|\nabla g(a)|^2/2\)。可逆性要求 \(I+tD^2g(a)\) 不退化。这与 Burgers 方程的分母具有同一几何含义，但高维中须检查整个矩阵。对于标量守恒律，下一节将用弱解与熵条件研究特征线交叉之后的延续。

\subsection{守恒律的位置}
\label{b2:subsec:070405}

设 \(d\geq1\) 为整数，\(F\in C^1(\mathbb R;\mathbb R^d)\)。形如 \(u_t+\operatorname{div}_xF(u)=0\) 的方程称为\term[shouhenglv]{守恒律}[conservation law]。若在时空开集 \(\Omega\subset\mathbb R^d\times\mathbb R\) 上 \(u\) 为实值 \(C^1\) 函数，则由链式法则它等价于拟线性方程 \(u_t+F'(u)\cdot\nabla u=0\)；不光滑时，两种写法的意义不同，后一项中的乘积未必已有定义。守恒形式可直接作用于测试函数，通过积分表达局部通量平衡。

\ProseParagraphBreak

线性连续性方程的速度由外部给定，密度与体积的补偿由 Jacobi 公式控制；非线性守恒律的速度随状态变化，投影交叉会使经典表示失效。两者都使用特征线，却不能把线性流的全局微分同胚结论原样移到非线性问题。本节建立了光滑阶段的局部构造和一个明确的失效机制，为进一步学习守恒律与 Hamilton--Jacobi 方程作准备。
